\chapter{Occupational Therapy}

\medskip
\noindent\textit{\textbf{Under review.} This chapter is an AI-assisted educational draft; it is not a substitute for professional medical advice.}

\section*{Definition and Overview}
Occupational therapy helps people participate in everyday activities, including self-care, work, education, rest, and recreation, after illness, injury, disability, developmental difference, or environmental change. The focus is on meaningful goals and practical function.

\section*{Assessment}
An occupational therapist assesses activities, cognition, movement, sensation, fatigue, communication, routines, environment, equipment, and the person's priorities. Assessment is collaborative and may involve family, carers, schools, employers, or other clinicians.

\section*{Interventions}
Treatment may include graded activity, strengthening and coordination, energy conservation, cognitive strategies, sensory or developmental support, splinting, assistive equipment, home modification, falls prevention, and return-to-work or school planning.

\section*{When to Seek Medical Care}
Occupational therapy does not replace urgent medical assessment. Call 000 for collapse, severe breathing difficulty, chest pain, stroke symptoms, or major injury. Seek routine review when function, safety, independence, or participation is declining.

\section*{Outcomes}
Progress is measured against agreed goals such as safer transfers, dressing, cooking, communication, school participation, work, or community access. Plans should be reviewed as abilities and circumstances change.

\section*{Prevention and Self-Care}
Use prescribed equipment correctly, practise agreed strategies, pace activities, reduce falls hazards, and tell the therapist about pain, skin problems, fatigue, or changes in function.

\section*{Expanded clinical context}
Occupational therapy is a clinical procedure, service, or rehabilitation topic. Interpretation should account for age, baseline health, medicines, comorbidities, goals, and access to care. This chapter supports discussion with a qualified clinician and does not establish a diagnosis.

\section*{Assessment and decision points}
Clarify the indication, expected benefit, alternatives, consent, preparation, recovery, and complications. Review history, allergies, medicines, pregnancy, previous reactions, and support needs. Document the main concern, time course, functional impact, relevant risks, and red flags. Use targeted investigations when their result is likely to change management.

\section*{Management and follow-up}
Safe care includes identity and site checks, appropriate analgesia or anaesthesia, monitoring, written discharge advice, rehabilitation, and planned follow-up. Explain expected improvement, when reassessment is needed, and how follow-up can be accessed. Avoid starting, stopping, or sharing prescription medicines without professional advice.

\section*{Safety-netting}
Seek urgent review for uncontrolled bleeding, severe pain, fever, spreading redness, wound separation, new neurological symptoms, breathing difficulty, or collapse.

\section*{Practical prevention}
Ask who to contact after hours and which medicines or activities should be paused or resumed. Keep written instructions and seek review if the topic recurs, changes, or does not improve as expected.

\section*{Limitations}
This educational expansion should be checked against current local guidance and authoritative topic-specific sources.
\clearpage
